\begin{mvbox}
设 $\mathbf{X}_{1},\mathbf{X}_{2},\cdots,\mathbf{X}_{n}$ 是来自均值为 $\boldsymbol{\mu}$ 与有限 (非奇异) 协方差 $\boldsymbol{\Sigma}$ 的总体的独立观测结果, 则对于大的 $n-p$,
\[
\sqrt{n}(\bar{\mathbf{X}}-\boldsymbol{\mu})\ \text{是近似的}\ N_{p}(\mathbf{0},\boldsymbol{\Sigma})\ \text{分布}
\]
且
\phantomsection\label{eq:4-28}
\begin{equation}
n(\bar{\mathbf{X}}-\boldsymbol{\mu})'\mathbf{S}^{-1}(\bar{\mathbf{X}}-\boldsymbol{\mu})\ \text{是近似的}\ \chi_{p}^{2}\ \text{分布} \tag{4-28}
\end{equation}
\end{mvbox}

在下面三节里, 我们考虑检验正态性假定的方法, 以及将非正态观测结果变换成近似正态的观测结果的方法.

\section{4.6 评估正态性假定}
\phantomsection\label{sec:4.6}

正如我们所指出的, 以后各章讨论的大多数统计方法都假定每个观测向量 $\mathbf{X}_{j}$ 来自多元正态分布. 另一方面, 在样本容量很大, 且所用的方法只取决于 $\mathbf{X}$ 的特性或形为 $n(\bar{\mathbf{X}}-\boldsymbol{\mu})'\mathbf{S}^{-1}(\bar{\mathbf{X}}-\boldsymbol{\mu})$ 的包含 $\bar{\mathbf{X}}$ 的距离时, 对于单个的观测结果的正态性假定并不是决定性的. 但在一定程度上, 由这些方法所作出的推断的质量随真实总体类似于多元正态形式的程度而定. 因此, 当数据显示出与多元正态的中度至极端的背离时, 必须要有能发现这种情况的方法存在.

我们要回答这个问题: 观测结果 $\mathbf{X}_{j}$ 是否违背它们来自正态总体的假定? 基于正态分布的性质, 我们已知正态变量的所有线性组合是正态的且多元正态密度的轮廓线是椭球面. 因此, 我们提出下述问题:

1.\quad $\mathbf{X}$ 的元素的边缘分布是否呈正态?

分量 $X_{i}$ 的几个线性组合是否呈正态?

2.\quad 根据各种特征的观测结果对作出的散布图, 是否给出正态总体所期望的椭圆形状?

3.\quad 是否存在应该进行检验以确保精确度的“杂乱”观测值?

我们的正态性研究将集中于对一维或二维中观测结果的特性 (例如, 边缘分布和散布图), 这将会很清楚. 不出所料, 结果表明要构造一个“好的”高于二维的联合正态性的全面的检验是困难的, 因为大量的东西可能出错. 在某种程度上, 我们必须为集中于一元和二元正态性的检验付出代价: 我们从来不能确保, 我们没有漏掉仅在高维中才能显示的某种特性. (这是可能的, 例如, 用正态的边缘分布构造一个非正态二元分布 [参见练习 4.8].) 然而非正态的许多类型通常是反映在边缘分布和散布图中. 另外, 对于大多数实际工作, 一维和二维的研究通常是足够的. 幸而, 在低维情况为正态而在高维情况为非正态的病态数据集, 在实际中并不是经常碰到.

\mvsub{评估一元边缘分布的正态性}

$n$ 较小时的点图和 $n>25$ 时的直方图有助于揭示一元分布的一个尾部比另一个长得多的的情况. 如果关于一个变量 $X_{i}$ 的直方图似乎较为对称, 我们可以通过计算在某个区间上的观测值的数目作进一步的检验. 一个一元正态分布属于区间 $(\mu_{i}-\sqrt{\sigma_{ii}},\mu_{i}+\sqrt{\sigma_{ii}})$ 的概率赋值为 0.683, 属于区间 $(\mu_{i}-2\sqrt{\sigma_{ii}},\mu_{i}+2\sqrt{\sigma_{ii}})$ 的概率为 0.954. 因此, 用一个大样本容量 $n$, 我们期望观察到位于区间 $(\bar{x}_{i}-\sqrt{s_{ii}},\bar{x}_{i}+\sqrt{s_{ii}})$ 的观测结果的比例 $\hat{p}_{1}$ 大约是 0.683. 类似的, 位于 $(\bar{x}_{i}-2\sqrt{s_{ii}},\bar{x}_{i}+2\sqrt{s_{ii}})$ 中的观测结果的比例 $\hat{p}_{2}$ 大约是 0.954. 利用关于 $\hat{p}_{i}$ 的抽样分布的正态近似 (参见文献 [9]), 我们观察到
